\chapter{Implementacja systemu}
\label{cha:pierwszyDokument}

\section{Architektura logiczna}
\label{sec:architektura_logiczna}

Projekt został podzielony na trzy warstwy:
\begin{itemize}
    \item \texttt{src/data} -- przygotowanie danych i podziały zbioru,
    \item \texttt{src/model} -- uczenie i ewaluacja klasyfikatorów,
    \item \texttt{src/inference} -- predykcja offline i online (kamera).
\end{itemize}

Warstwa wejściowa CLI znajduje się w \texttt{main.py} i deleguje logikę do
modułów w \texttt{src/}. Takie podejście upraszcza testowanie i ogranicza
sprzężenie między interfejsem użytkownika a logiką przetwarzania.

\section{Przetwarzanie danych}
\label{sec:przetwarzanie_danych}

Klasa \texttt{CsvDataset} ładuje plik CSV i wymusza spójność schematu
(obecność kolumny \texttt{label} oraz dokładnie 784 cechy obrazu). Dla każdej
próbki zwracana jest para:
\begin{center}
    \texttt{(feature\_vector: np.ndarray, label: str)}
\end{center}

Wektor cech jest typu \texttt{float32} i przyjmuje wartości w $[0,1]$.
Etykiety numeryczne przekształcane są do liter poprzez mapowanie:
\begin{equation}
    y_{letter} = \mathrm{chr}(\mathrm{ord}('A') + y_{int})
\end{equation}

Budowa macierzy cech realizowana jest przez \texttt{build\_feature\_matrix},
a podział danych przez \texttt{stratified\_split}. Domyślnie stosowany jest
podział 80/10/10 z zachowaniem proporcji klas.

\section{Uczenie modeli}
\label{sec:uczenie_modeli}

Moduł \texttt{src/model/train.py} obsługuje trzy warianty klasyfikatorów:
\begin{itemize}
    \item \textbf{SVC (RBF)} -- domyślny klasyfikator oparty o jądro radialne,
    \item \textbf{Random Forest} -- las losowy o 200 drzewach,
    \item \textbf{CNN (Keras)} -- sieć konwolucyjna dla wejścia 28$\times$28.
\end{itemize}

Modele klasyczne serializowane są przez \texttt{joblib} do formatu
\texttt{.pkl}. Model konwolucyjny zapisywany jest jako \texttt{.keras}
oraz osobny plik mapowania etykiet \texttt{.labels.json}.

\section{Ewaluacja jakości}
\label{sec:ewaluacja_jakosci}

Funkcja \texttt{evaluate} oblicza dokładność globalną oraz raport klasyfikacji
\texttt{classification\_report} biblioteki scikit-learn \cite{sklearn_docs}.
Wyniki zwracane są również jako słownik, co upraszcza dalsze automatyczne
raportowanie.

\section{Inferencja: CSV i kamera}
\label{sec:inferencja_csv_kamera}

Inferencja offline z jednego wiersza CSV realizowana jest przez
\texttt{CsvRowSource}. Inferencja online wykorzystuje klasę \texttt{CameraSource},
która:
\begin{enumerate}
    \item pobiera klatkę z kamery (OpenCV \cite{opencv_docs}),
    \item wykrywa dłoń przez MediaPipe GestureRecognizer \cite{mediapipe_docs},
    \item wyznacza obszar dłoni i dokonuje kadrowania,
    \item przeskalowuje obraz do 28$\times$28 oraz normalizuje piksele,
    \item przekazuje wektor 784 cech do modelu.
\end{enumerate}

W przypadku braku detekcji dłoni moduł zwraca wektor zerowy, dzięki czemu
interfejs inferencji pozostaje stabilny.

\section{Interfejs CLI}
\label{sec:interfejs_cli}

Plik \texttt{main.py} udostępnia trzy tryby:
\begin{itemize}
    \item \texttt{train} -- trening wybranego klasyfikatora,
    \item \texttt{eval} -- ewaluacja zapisanego modelu,
    \item \texttt{infer} -- predykcja z CSV lub kamery.
\end{itemize}

Przykładowe użycie:
\begin{lstlisting}[style=stylePython,caption={Przyklady uruchomien CLI}]
uv run python main.py train --classifier svc --model artifacts/svc_model.pkl
uv run python main.py eval --model artifacts/svc_model.pkl --dataset data/raw/sign_mnist_test.csv
uv run python main.py infer --source csv --row 42 --dataset data/raw/sign_mnist_test.csv
uv run python main.py infer --source camera --camera 0 --model artifacts/svc_model.pkl
\end{lstlisting}
